\section{Profils}

\subsection{Exercice 4.1 -- Profil de ressources}
\textbf{Objectif.} Vérifier \code{RESOURCE\_LIMIT}, créer le profil \code{P\_RH\_TEST}
(\code{SESSIONS\_PER\_USER 2}, \code{IDLE\_TIME 1}) et l'attribuer à LINA, puis observer
les limites de sessions simultanées et d'inactivité.

\shots{4_1_setup}
\shots{4_1_a}
\shots{4_1_b}
\shots{4_1_third}
\shots{4_1_idle}
\shots{4_1_idle_status}

\textbf{Observations et réponses.}
\begin{enumerate}
  \item \code{RESOURCE\_LIMIT} vaut \code{TRUE} : les limites de profil sont appliquées.
  \item Le profil \code{P\_RH\_TEST} est créé avec \code{SESSIONS\_PER\_USER = 2} et
        \code{IDLE\_TIME = 1} (1 minute), puis attribué à LINA.
  \item La troisième session LINA renvoie \code{ORA-02391: exceeded simultaneous
        SESSIONS\_PER\_USER limit} : le plafond de 2 sessions simultanées est respecté.
  \item Après \code{HOST sleep 150} (> IDLE\_TIME 1 minute), la requête renvoie
        \code{ORA-02396: exceeded maximum idle time, please connect again}. Pendant
        l'attente, le DBA voit la session LINA en statut \code{ACTIVE} dans \code{V\$SESSION},
        puis elle disparaît (\emph{sniped} par Oracle).
  \item Dans \code{DBA\_PROFILES}, la valeur \code{DEFAULT} dans la colonne \code{LIMIT}
        signifie que le paramètre reprend la valeur définie dans le profil \code{DEFAULT}.
\end{enumerate}
\ecart{Aucun. Le comportement est conforme à l'énoncé.}

\subsection{Exercice 4.2 -- Verrouillage après échecs}
\textbf{Objectif.} Configurer le verrouillage de compte après échecs de connexion
(\code{FAILED\_LOGIN\_ATTEMPTS 3}, \code{PASSWORD\_LOCK\_TIME 2/1440}) et en observer
l'effet ainsi que le déverrouillage automatique.

\shots{4_2}

\textbf{Observations et réponses.}
\begin{enumerate}
  \item Les trois premières connexions LINA avec un mauvais mot de passe renvoient
        \code{ORA-01017}. La quatrième, avec le bon mot de passe, renvoie
        \code{ORA-28000: The account is locked; login denied.} Le compte est verrouillé.
  \item \code{DBA\_USERS} affiche \code{ACCOUNT\_STATUS = LOCKED(TIMED)} avec un
        \code{LOCK\_DATE} renseigné.
  \item Après une attente de 130 secondes (> 2 minutes de \code{PASSWORD\_LOCK\_TIME}),
        la reconnexion de LINA réussit : le déverrouillage est automatique.
\end{enumerate}
\ecart{Aucun. Conforme à l'énoncé.}

\subsection{Exercice 4.3 -- Expiration, grâce et réutilisation}
\textbf{Objectif.} Configurer la réutilisation des mots de passe
(\code{PASSWORD\_REUSE\_TIME 30}, \code{PASSWORD\_REUSE\_MAX 3}), puis tester l'expiration
et la période de grâce avec \code{PASSWORD\_LIFE\_TIME} et \code{PASSWORD\_GRACE\_TIME}
réglés à 1 minute.

\shots{4_3}

\textbf{Observations et réponses.}
\begin{enumerate}
  \item Les règles \code{PASSWORD\_REUSE\_TIME = 30} et \code{PASSWORD\_REUSE\_MAX = 3}
        sont appliquées à \code{P\_RH\_TEST}.
  \item LINA change son mot de passe, puis tente de reprendre l'ancien immédiatement :
        la reprise a \textbf{réussi} (pas de \code{ORA-28007}) : l'historique de
        réutilisation n'a pas bloqué la réutilisation immédiate dans ce contexte.
        \ecart{L'énoncé 19c attend un refus (ORA-28007) en réutilisant l'ancien mot de
        passe ; en 26ai la reprise a fonctionné. Cela est consigné tel quel dans le log.}
  \item Avec \code{PASSWORD\_LIFE\_TIME = 1/1440} et \code{PASSWORD\_GRACE\_TIME =
        1/1440} : après 1 minute, la connexion est acceptée avec l'avertissement
        \code{ORA-28098} et \code{ACCOUNT\_STATUS = EXPIRED(GRACE)}. Après 2 minutes
        supplémentaires, la connexion renvoie \code{ORA-28001} et le compte passe en
        \code{EXPIRED}.
  \item Restauration à \code{PASSWORD\_LIFE\_TIME 90} / \code{PASSWORD\_GRACE\_TIME 7} :
        LINA reçoit un mot de passe neuf (\code{Lina\#2032}) et repasse en \code{OPEN}.
\end{enumerate}

\subsection{Exercice 4.4 -- Complexité des mots de passe}
\textbf{Objectif.} Associer une fonction de vérification de complexité
(\code{ORA12C\_VERIFY\_FUNCTION}) au profil et observer le rejet des mots de passe faibles
ainsi que l'obligation de la clause \code{REPLACE}.

\shots{4_4}

\textbf{Observations et réponses.}
\begin{enumerate}
  \item \code{ORA12C\_VERIFY\_FUNCTION} existe bien (objet \code{FUNCTION} propriétaire
        \code{SYS}). Elle est associée à \code{P\_RH\_TEST}.
  \item Tentatives de mots de passe faibles (\code{lina}, \code{abcdefgh},
        \code{Lina2026x}) sans clause \code{REPLACE} : toutes renvoient
        \code{ORA-28221: REPLACE clause not specified}. Quand une fonction de vérification
        est définie, Oracle exige la clause \code{REPLACE} avant d'évaluer la complexité.
  \item Le mot de passe \code{Lina\#2033} (contient le nom d'utilisateur) renvoie
        \code{ORA-28003: password contains the user name}, même avec la clause
        \code{REPLACE}.
\end{enumerate}
\ecart{L'énoncé 19c attend \code{ORA-28003} pour les mots de passe faibles. En 26ai, sans
clause \code{REPLACE}, c'est \code{ORA-28221} qui est renvoyé en premier (exigence de la
clause REPLACE liée à la fonction de vérification). Le mot de passe de LINA reste
\code{Lina\#2032} (compte \code{OPEN}).}

\subsection{Exercice 4.5 -- Profil de production et synthèse}
\textbf{Objectif.} Créer un profil de production \code{P\_RH\_STANDARD}, le déployer sur
LINA, OMAR et SARA, supprimer \code{P\_RH\_TEST} avec \code{CASCADE} et observer le sort
des utilisateurs concernés.

\shots{4_5}

\textbf{Observations et réponses.}
\begin{enumerate}
  \item Le profil \code{P\_RH\_STANDARD} est créé avec des limites de production
        (\code{SESSIONS\_PER\_USER 2}, \code{IDLE\_TIME 30}, \code{FAILED\_LOGIN\_ATTEMPTS
        3}, \code{PASSWORD\_LIFE\_TIME 90}, \code{PASSWORD\_REUSE\_TIME 365},
        \code{PASSWORD\_VERIFY\_FUNCTION ORA12C\_VERIFY\_FUNCTION}, etc.).
  \item \code{DROP PROFILE p\_rh\_test CASCADE} supprime le profil et réaffecte
        automatiquement LINA, OMAR et SARA au profil \code{DEFAULT} (visible dans
        \code{DBA\_USERS}). Un utilisateur dont le profil est supprimé avec
        \code{CASCADE} retombe donc sur le profil \code{DEFAULT}.
  \item Les trois utilisateurs sont ensuite passés sur \code{P\_RH\_STANDARD} ; leur
        \code{ACCOUNT\_STATUS} reste \code{OPEN}.
\end{enumerate}
\ecart{Aucun. Conforme à l'énoncé.}
